%% =====================================================================
%%  T-06-03  The Third-Party Lie Rule
%%  -------------------------------------------------------------------
%%  One idea per file. Core is mandatory. Slots in canonical order.
%%  Max 700 words. No layout commands. No filler. New source material
%%  for this entry goes in sources/<BOOK>/T-06-03.tex -- never by
%%  rewriting this file (docs/RULES.md R0.1).
%% =====================================================================
%% @kind: topic
%% @id: T-06-03
%% @title: The Third-Party Lie Rule
%% @subtitle: How they treat someone else tells you how they will treat you
%% @chapter: c06
%% @order: 30
%% @status: stable
%% @sources: B1
%% @updated: 2026-09-19

\topic{T-06-03}{The Third-Party Lie Rule}{How they treat someone else tells you how they will treat you}

\begin{core}
Watch a person lie to, or cheat, or dismiss someone who is not you. Their conduct toward third parties is not a separate department of their character; it is the same rule they will apply to you once your classification changes.
\end{core}
\begin{mechanism}
People operate from an internal code that determines whom it is acceptable to deceive. The code is rarely organised around persons; it is organised around categories --- people who matter, people who do not, people who can retaliate, people who cannot. You are currently in a favourable category, which is why you are seeing the good behaviour rather than the code. Categories are reassignable at low cost, and the reassignment happens the moment you stop being useful or start being inconvenient. Observing a third party being mistreated is therefore a direct observation of the rule that will govern you later, taken at the only time it is visible.
\end{mechanism}
\begin{tells}
  \item They lie to someone in front of you, and then assure you that you are different.
  \item Service staff, subordinates and strangers are treated noticeably worse than people who can affect them.
  \item An absent third party is described in terms that make exploiting them sound reasonable.
  \item Loyalty is offered as a special exemption from a general practice.
  \item They recount past deceptions of others with approval rather than with discomfort.
\end{tells}
\begin{moves}
  \item Establish a category boundary with the target --- friend, outsider, adversary --- and behave by category rather than by person.
  \item Offer the counterpart a privileged position inside the category system.
  \item Demonstrate the practice on someone else first, where the cost of demonstration is low.
  \item Explain the practice as a response to how the world works rather than as a choice.
\end{moves}
\begin{counters}
  \item Treat observed conduct toward third parties as data about you, on a delay.
  \item Reject special-exemption offers. A claim that you would be treated differently, from a person who treats others badly, is a statement about your current category rather than about their character.
  \item Ask how they describe people they have fallen out with. The description is a rehearsal.
  \item Assume your category will change, and structure commitments so that it does not matter --- written terms, no dependence, an exit that does not require their cooperation.
\end{counters}
\begin{limits}
Everyone is harsher with some people than others, and a single observation proves little; the rule reads a settled pattern, not an incident. It is also asymmetric in a way that matters --- being treated well by someone who treats others badly is genuinely pleasant and genuinely unstable, which makes the observation hard to act on. Finally the rule is unavailable in exactly the situation it is most needed: the operator will not mistreat a third party in front of a mark they are still working.
\end{limits}

\sources{B1}
\seesources
\seealso{T-03-02, T-06-01, T-06-04}
